\section{Experimental design}
\label{sec:experiments}
\subsection{Development protocol and landmark evaluation}
The 801 development patients were partitioned into five patient-level folds. Fold 1 used 640 training and 161 validation patients; folds 2--5 each used 641 training and 160 validation patients. The same 175-patient holdout was evaluated by every fold model and every variant. Thus the five test evaluations per variant are repeated evaluations of overlapping trained models on the same patients, rather than five independent test cohorts. Patient identities were disjoint between training, validation, and holdout within the current split.

At each training epoch, a patient contributed one randomly selected available T0, T1, or T2 prefix. Only MRI observations available at or before the selected landmark were classifier inputs. Later observed visits, including T3, could supply latent targets for the auxiliary loss on training patients. Missing targets were masked. The 23 patients with only T0 remained eligible for pCR training but supplied no future-state loss. At evaluation, all three landmarks were evaluated for the 175 holdout patients, each of whom had T0, T1, and T2. T3 was not an evaluation landmark and was not used as observed input for these predictions.

\subsection{Model comparisons}
The principal comparison comprises three variants trained with seeds 21 and 44 and the same five development folds. \nojepa{} uses direct pCR readout and a zero future-state loss weight. \direct{} uses the same direct readout with a future-state loss weight of 0.5. \fullmodel{} retains that auxiliary weight and replaces the direct readout with a concatenation of direct pCR features, attention-pooled predicted terminal state, and attention-pooled predicted change. All variants retain the shared image encoder, clinical inputs, assigned-arm encoding, fusion architecture, and training schedule.

The \direct{}--\nojepa{} contrast most directly tests adding future-state supervision. The \fullmodel{}--\direct{} contrast jointly changes the use of forecast features, spatial pooling, and classifier input width. It cannot isolate attention alone. For that purpose, the supplementary \meanmodel{} comparison uses forecast readout with mean spatial pooling and otherwise matches the forecast-attention configuration. All manuscript-facing results use the same equal aggregation of seeds 21 and 44. No VQ-latent reconstruction-head, change-weighted-loss, or state-conditioned-head variants are included.

\subsection{Implementation and optimization}
The network uses a model width of 768, eight layers in each of the spatial encoder, multimodal fusion module, and predictor, 16 attention heads with 64 dimensions per head, and feed-forward width 3,072. The 1,024-dimensional concatenated attention-head representation is projected back to the 768-dimensional residual stream. The spatial token grid is $2\times4\times4$. The configuration uses eight state queries, four pCR queries, dropout 0.3, and an EMA coefficient of 0.999. The pretrained convolutional encoder is fine-tuned jointly with the remaining network. Its training-patient provenance requires additional documentation, as discussed in Section~\ref{sec:limitations}.

We use AdamW \citep{adamw2019} with an initial base learning rate of $10^{-4}$ for both the encoder and the main network, weight decay 0.05, batch size 16, no gradient accumulation, and gradient-norm clipping at 1.0. Training uses BF16 mixed precision. A ten-epoch warm-up increases the learning rate from $10^{-5}$ to $10^{-4}$, followed by cosine decay to $10^{-6}$ at epoch 100. Training is capped at 200 epochs, with early stopping after 50 epochs without improvement. The best checkpoint is selected independently in each fold using validation patient-equal-prefix F1 at threshold 0.5. This is an F1 selection protocol, even though AUROC is a principal reporting measure. Deterministic execution was not enabled, so a fixed seed does not imply bitwise reproducibility across hardware or software changes.

\subsection{Metrics, aggregation, and patient-level uncertainty}
We report 14 complementary metrics: AUROC, non-interpolated average precision (AP), F1, accuracy, sensitivity, specificity, positive and negative predictive value (PPV/NPV), balanced accuracy, Matthews correlation coefficient (MCC), Youden index, negative log-likelihood (NLL), Brier score, and expected calibration error (ECE) with ten equal-width probability bins. Confusion-matrix counts are additionally retained. Thresholded primary metrics use 0.5. MCC and balanced accuracy help interpret performance under class imbalance~\citep{chicco2020mcc}; positive prevalence is $56/175=0.32$. No holdout threshold optimization or probability recalibration was performed.

Let $M$ denote a metric, $p_{isf\ell}$ the probability for patient $i$, training seed $s$, fold $f$, and landmark $\ell$, and $y_i$ the pCR label. Our single reported landmark summary is
\begin{equation}
\begin{aligned}
 \overline M_{s\ell}&=\frac{1}{5}\sum_{f=1}^{5}M\bigl(\{y_i,p_{isf\ell}\}_{i=1}^{175}\bigr),\\
 \overline M_{\ell}&=\frac{\overline M_{21,\ell}+\overline M_{44,\ell}}{2},\\
 \overline M_{\mathrm{macro}}&=\frac{1}{3}\sum_{\ell=0}^{2}\overline M_{\ell}.
\end{aligned}
\end{equation}
We do not average probabilities across folds or seeds before computing $M$. Each metric is calculated independently in each fitted model; within-seed five-fold means are then equally averaged across seeds. A mean F1 across landmarks is likewise not recomputed after pooling landmark rows. Only the combined two-seed estimate is presented in manuscript figures and tables.

For exploratory uncertainty estimates, we draw 2,000 bootstrap samples of the 175 patients with replacement. Every patient draw is shared across both seeds, all fold models, variants, and landmarks, preserving pairing and repeated observations. Within each resample, metrics are computed separately per fold and then aggregated exactly as above. Percentile 2.5th and 97.5th quantiles form the intervals. They condition on the fitted models and selected analysis; they do not capture retraining, seed-selection, or model-selection uncertainty. Neither ten fitted models nor three repeated landmarks are treated as independent patient samples. We do not make multiplicity-adjusted confirmatory significance claims.

ROC curves are interpolated on a common false-positive-rate grid; PR curves are summarized on a recall grid. Curves and their patient-bootstrap bands follow the same fold-then-seed averaging. Separately calculated mean AUROC/AP values are reported alongside the curves, since interpolation and curve averaging can change the displayed area. Reliability diagrams use five equal-count probability-rank groups within each fitted model, then average corresponding coordinates; these descriptive bins are distinct from the ten fixed bins used for ECE.

\subsection{Operating-point and hypothetical decision analyses}
As a threshold sensitivity analysis, Youden index is maximized on each fold's validation patients separately at T0, T1, and T2. The highest threshold breaks exact ties. That threshold is then fixed for the corresponding holdout evaluation; no holdout prediction chooses it. Bootstrap intervals condition on the fitted threshold rather than refitting it in each resample.

We additionally compute exploratory decision curves~\citep{vickers2006dca}. At risk threshold $q$, the generic positive action is applied when predicted pCR probability is at least $q$. Its mathematical net benefit is
\begin{equation}
 \mathrm{NB}(q)=\frac{\mathrm{TP}(q)}{n}
 -\frac{\mathrm{FP}(q)}{n}\frac{q}{1-q}.
\end{equation}
We calculate this separately for every fitted model over thresholds 0.05--0.80, then apply fold-then-seed aggregation. Treat-all and treat-none are references. The calculation uses uncalibrated probabilities and does not validate a particular escalation, de-escalation, or other management action.

\subsection{Exploratory subgroup and mechanism analyses}
We examine released HR/HER2 combinations, age below or at least 50 years, native source namespace, and saved predicted-component crop coverage. Age 50 is an exploratory reporting cutpoint. Subgroup AUROC/AP are suppressed when fewer than ten patients or fewer than five outcomes in either class are available. Paired within-subgroup patient resampling preserves all seed/fold/landmark pairing; single-class draws are omitted for AUROC intervals and valid-draw counts retained. Intervals are unadjusted, and no subgroup-interaction hypothesis is tested. Source namespaces share the same clinical release and are not independent external cohorts.

Mechanism analysis is restricted to archived fixed-training-probe diagnostics collected every five epochs. These probes sample up to 32 training patients and are not full validation or holdout forecast evaluations. Unlike the missing-target-masked training objective, the archived forecast diagnostic includes only continuous observed visit chains: the real landmark, target, and every intermediate canonical visit must be available. A gap excludes that landmark--target pair, and eligible-patient counts therefore vary by horizon. We report a descriptive common epoch 20 and common-epoch trajectories. Forecast skill is $1-\mathrm{L1}_{\mathrm{forecast}}/\mathrm{L1}_{\mathrm{copy}}$ within the same snapshot's EMA representation; the copy forecast holds the observed landmark state constant. This normalization avoids comparing raw latent errors between separately learned teacher coordinates. Predictor gradient norms from the pCR and weighted rollout objectives are descriptive checks of gradient routing, not independent statistical samples or clinical effects.

\subsection{Scope of the selected-seed-pair analysis}
The choice to average seeds 21 and 44 was made retrospectively after multiple-seed results were available, including holdout results. It was not a prospectively registered seed policy. Using the same pair and weights for every model preserves a matched comparison, but does not remove selection bias or fully characterize training randomness. The analysis is exploratory. A confirmatory study should lock model selection and the seed policy before evaluating independent patients.
